

%    _   _       _               _   
%   /_\ | |__ __| |_ _ _ __ _ __| |_ 
%  / _ \| '_ (_-<  _| '_/ _` / _|  _|
% /_/ \_\_.__/__/\__|_| \__,_\__|\__|
%

\begin{abstract} 
\noindent Representation theory is an immense topic. Most sources we have
encountered emphasize rigor and depth over conceptual clarity, especially for
physics. Such presentations tend to conceal the exceptional regularity and
beauty of the subject just due to their length. Here, we present the bare
minimum, from the physicist's perspective, following parts of
Schwichtenberg\cite{schwichtenberg2018} closely. In addition to resolving
ambiguities as we find them, we specifically bridge classical physics, where
representation theory is not often invoked, but where we find it quite
illuminating, and quantum physics, where the physics student typically first
encounters the subject. We consider this unfortunate, because the student both
misses the clarifying utility of representation theory in classical physics and
because the student of quantum physics already has many steep hills to climb.
We deem it better to cover representation theory as a separate topic pertinent
to both realms of physics.
% We skip elaboration of his
% (good!) examples to bring definitions textually closer to their usages and to
% clarify some points. We invoke other sources as needed. 
\end{abstract}

